% ============================================================
% CHAPTER 4: DATABASE DESIGN
% ============================================================
\chapter{Database Design}

This chapter presents the database design for the SafeGuard system, including the database overview, Entity-Relationship modelling, detailed schema definitions, table structures with their attributes, key constraints, relationships, and a comprehensive data dictionary.

\section{Database Overview}

The SafeGuard system uses \textbf{MySQL 8.0} as its relational database management system (RDBMS). MySQL was selected for its reliability, performance, widespread adoption in web applications, and excellent compatibility with Python's SQLAlchemy ORM through the PyMySQL driver.

The database, named \texttt{safety\_alert\_db}, is managed through SQLAlchemy's declarative base system, which maps Python classes to database tables. The schema is automatically created at application startup using \texttt{Base.metadata.create\_all()}, eliminating the need for manual SQL migration scripts during development.

The database follows the \textbf{Third Normal Form (3NF)} of normalization, ensuring minimal data redundancy and optimal data integrity. The schema comprises five primary tables, each representing a distinct entity in the system's domain model.

\begin{table}[H]
  \centering
  \caption{Database Tables Overview}
  \label{tab:db-overview}
  \begin{tabularx}{\textwidth}{@{}l X c@{}}
    \toprule
    \textbf{Table Name} & \textbf{Purpose} & \textbf{Est. Records} \\
    \midrule
    \texttt{users} & Stores user account information, credentials, and profile data & 1--1000 \\
    \texttt{emergency\_contacts} & Stores user-managed emergency contact details & 1--50 per user \\
    \texttt{sos\_alerts} & Records all SOS emergency alert events & 0--100 per user \\
    \texttt{incident\_reports} & Stores structured incident reports with status tracking & 0--50 per user \\
    \texttt{activity\_logs} & Logs all significant user actions for audit purposes & 5--200 per user \\
    \bottomrule
  \end{tabularx}
\end{table}

\section{Entity-Relationship Diagram}

Figure~\ref{fig:er-diagram} presents the Entity-Relationship (ER) diagram for the SafeGuard database, illustrating the entities, attributes, and relationships between tables.

\begin{figure}[H]
  \centering
  \begin{tikzpicture}[
    entity/.style={rectangle, draw, minimum width=2.5cm, minimum height=0.7cm, text centered, font=\small\bfseries, fill=blue!10},
    weakentity/.style={rectangle, draw, dashed, minimum width=2.5cm, minimum height=0.7cm, text centered, font=\small\bfseries, fill=red!5},
    relation/.style={diamond, draw, minimum width=1.8cm, minimum height=0.6cm, text centered, font=\tiny},
    pkattr/.style={ellipse, draw, fill=yellow!15, font=\tiny},
    fkattr/.style={ellipse, draw, fill=green!10, font=\tiny},
    normattr/.style={ellipse, draw, font=\tiny},
    line/.style={thick},
  ]

    % User Entity
    \node[entity] (user) at (0, 0) {users};
    \node[pkattr] (uid) at (-3, 1) {id (PK)};
    \node[normattr] (uname) at (-3, 0) {first\_name};
    \node[normattr] (uemail) at (-3, -1) {email};
    \node[normattr] (upass) at (-3, -2) {password\_hash};
    \node[normattr] (urole) at (3, 1) {role};
    \node[normattr] (uphone) at (3, 0) {phone};
    \node[normattr] (uactive) at (3, -1) {is\_active};
    \draw[line] (user) -- (uid);
    \draw[line] (user) -- (uname);
    \draw[line] (user) -- (uemail);
    \draw[line] (user) -- (upass);
    \draw[line] (user) -- (urole);
    \draw[line] (user) -- (uphone);
    \draw[line] (user) -- (uactive);

    % Emergency Contact Entity
    \node[entity] (contact) at (-6, -5) {emergency\_contacts};
    \node[pkattr] (cid) at (-8, -4) {id (PK)};
    \node[fkattr] (cuid) at (-8, -5) {user\_id (FK)};
    \node[normattr] (cname) at (-8, -6) {name};
    \node[normattr] (cphone) at (-4, -4) {phone};
    \node[normattr] (cpriority) at (-4, -5) {priority};
    \draw[line] (contact) -- (cid);
    \draw[line] (contact) -- (cuid);
    \draw[line] (contact) -- (cname);
    \draw[line] (contact) -- (cphone);
    \draw[line] (contact) -- (cpriority);

    % SOS Alert Entity
    \node[entity, fill=red!8] (sos) at (0, -5) {sos\_alerts};
    \node[pkattr] (sid) at (-1.5, -4) {id (PK)};
    \node[fkattr] (suid) at (1.5, -4) {user\_id (FK)};
    \node[normattr] (slat) at (-1.5, -6) {latitude};
    \node[normattr] (stype) at (1.5, -6) {alert\_type};
    \node[normattr] (sstatus) at (0, -6.8) {status};
    \draw[line] (sos) -- (sid);
    \draw[line] (sos) -- (suid);
    \draw[line] (sos) -- (slat);
    \draw[line] (sos) -- (stype);
    \draw[line] (sos) -- (sstatus);

    % Incident Entity
    \node[entity, fill=orange!8] (incident) at (6, -5) {incident\_reports};
    \node[pkattr] (iid) at (4, -4) {id (PK)};
    \node[fkattr] (iuid) at (8, -4) {user\_id (FK)};
    \node[normattr] (iincid) at (4, -6) {incident\_id};
    \node[normattr] (itype) at (8, -6) {incident\_type};
    \node[normattr] (isev) at (6, -6.8) {severity};
    \draw[line] (incident) -- (iid);
    \draw[line] (incident) -- (iuid);
    \draw[line] (incident) -- (iincid);
    \draw[line] (incident) -- (itype);
    \draw[line] (incident) -- (isev);

    % Activity Log Entity
    \node[entity, fill=purple!8] (activity) at (0, -10) {activity\_logs};
    \node[pkattr] (aid) at (-1.5, -9) {id (PK)};
    \node[fkattr] (auid) at (1.5, -9) {user\_id (FK)};
    \node[normattr] (aaction) at (0, -11) {action};
    \draw[line] (activity) -- (aid);
    \draw[line] (activity) -- (auid);
    \draw[line] (activity) -- (aaction);

    % Relations
    \node[relation] (r1) at (-3, -2.5) {has};
    \node[relation] (r2) at (0, -2.5) {creates};
    \node[relation] (r3) at (3, -2.5) {reports};
    \node[relation] (r4) at (0, -7.5) {generates};

    \draw[line] (user) -- (r1);
    \draw[line] (r1) -- (contact);
    \draw[line] (user) -- (r2);
    \draw[line] (r2) -- (sos);
    \draw[line] (user) -- (r3);
    \draw[line] (r3) -- (incident);
    \draw[line] (user) -- (r4);
    \draw[line] (r4) -- (activity);

  \end{tikzpicture}
  \caption{Entity-Relationship Diagram}
  \label{fig:er-diagram}
\end{figure}

\section{Database Schema}

The SafeGuard database schema is defined through SQLAlchemy ORM model classes. Each model class maps to a corresponding database table with columns, constraints, and relationships defined at the class level. The ORM approach ensures database portability and type-safe data access.

\section{Table Definitions}

\subsection{Users Table}

The \texttt{users} table is the central entity storing user account information and authentication credentials.

\begin{table}[H]
  \centering
  \caption{Users Table Schema}
  \label{tab:schema-users}
  \begin{tabularx}{\textwidth}{@{}l l l X l@{}}
    \toprule
    \textbf{Column} & \textbf{Type} & \textbf{Constraints} & \textbf{Description} & \textbf{Default} \\
    \midrule
    id & INTEGER & PK, AUTO\_INCREMENT & Unique user identifier & --- \\
    first\_name & VARCHAR(100) & NOT NULL & User's first name & --- \\
    last\_name & VARCHAR(100) & NOT NULL & User's last name & --- \\
    email & VARCHAR(255) & UNIQUE, NOT NULL, INDEX & User's email address & --- \\
    phone & VARCHAR(20) & NOT NULL & Contact phone number & --- \\
    password\_hash & VARCHAR(255) & NOT NULL & bcrypt hashed password & --- \\
    profile\_picture & VARCHAR(500) & NULLABLE & Path to profile image & NULL \\
    role & ENUM & NOT NULL & User or admin role & 'user' \\
    is\_active & BOOLEAN & NOT NULL & Account active status & TRUE \\
    created\_at & TIMESTAMP & NOT NULL & Account creation time & NOW() \\
    updated\_at & TIMESTAMP & NOT NULL & Last update timestamp & NOW() \\
    \bottomrule
  \end{tabularx}
\end{table}

\subsection{Emergency Contacts Table}

The \texttt{emergency\_contacts} table stores the emergency contact details for each user.

\begin{table}[H]
  \centering
  \caption{Emergency Contacts Table Schema}
  \label{tab:schema-contacts}
  \begin{tabularx}{\textwidth}{@{}l l l X l@{}}
    \toprule
    \textbf{Column} & \textbf{Type} & \textbf{Constraints} & \textbf{Description} & \textbf{Default} \\
    \midrule
    id & INTEGER & PK, AUTO\_INCREMENT & Unique contact identifier & --- \\
    user\_id & INTEGER & FK $\to$ users.id, NOT NULL, INDEX & Owner user reference & --- \\
    name & VARCHAR(200) & NOT NULL & Contact person's name & --- \\
    relationship & VARCHAR(100) & NOT NULL & Relationship type & --- \\
    phone & VARCHAR(20) & NOT NULL & Contact phone number & --- \\
    email & VARCHAR(255) & NULLABLE & Contact email address & NULL \\
    priority & ENUM & NOT NULL & Priority level & 'medium' \\
    created\_at & TIMESTAMP & NOT NULL & Creation timestamp & NOW() \\
    updated\_at & TIMESTAMP & NOT NULL & Last update timestamp & NOW() \\
    \bottomrule
  \end{tabularx}
\end{table}

\subsection{SOS Alerts Table}

The \texttt{sos\_alerts} table records all emergency alert events triggered by users.

\begin{table}[H]
  \centering
  \caption{SOS Alerts Table Schema}
  \label{tab:schema-sos}
  \begin{tabularx}{\textwidth}{@{}l l l X l@{}}
    \toprule
    \textbf{Column} & \textbf{Type} & \textbf{Constraints} & \textbf{Description} & \textbf{Default} \\
    \midrule
    id & INTEGER & PK, AUTO\_INCREMENT & Unique alert identifier & --- \\
    user\_id & INTEGER & FK $\to$ users.id, NOT NULL, INDEX & Alert creator reference & --- \\
    latitude & NUMERIC(10,8) & NULLABLE & GPS latitude coordinate & NULL \\
    longitude & NUMERIC(11,8) & NULLABLE & GPS longitude coordinate & NULL \\
    address & TEXT & NULLABLE & Textual location description & NULL \\
    message & TEXT & NULLABLE & Custom emergency message & NULL \\
    alert\_type & ENUM & NOT NULL & Alert classification & 'manual' \\
    status & ENUM & NOT NULL, INDEX & Alert lifecycle status & 'active' \\
    created\_at & TIMESTAMP & NOT NULL & Alert creation time & NOW() \\
    updated\_at & TIMESTAMP & NOT NULL & Last status update & NOW() \\
    \bottomrule
  \end{tabularx}
\end{table}

\subsection{Incident Reports Table}

The \texttt{incident\_reports} table stores structured safety incident reports.

\begin{table}[H]
  \centering
  \caption{Incident Reports Table Schema}
  \label{tab:schema-incident}
  \begin{tabularx}{\textwidth}{@{}l l l X l@{}}
    \toprule
    \textbf{Column} & \textbf{Type} & \textbf{Constraints} & \textbf{Description} & \textbf{Default} \\
    \midrule
    id & INTEGER & PK, AUTO\_INCREMENT & Internal record ID & --- \\
    user\_id & INTEGER & FK $\to$ users.id, NOT NULL, INDEX & Report creator reference & --- \\
    incident\_id & VARCHAR(50) & UNIQUE, NOT NULL, INDEX & Public incident ID (INC-XXXXXXXX) & --- \\
    incident\_type & ENUM & NOT NULL & Type of incident & --- \\
    description & TEXT & NOT NULL & Detailed incident description & --- \\
    location & VARCHAR(500) & NOT NULL & Location text & --- \\
    latitude & NUMERIC(10,8) & NULLABLE & GPS latitude & NULL \\
    longitude & NUMERIC(11,8) & NULLABLE & GPS longitude & NULL \\
    severity & ENUM & NOT NULL & Severity level & 'medium' \\
    image & VARCHAR(500) & NULLABLE & Evidence image path & NULL \\
    status & ENUM & NOT NULL & Investigation status & 'pending' \\
    created\_at & TIMESTAMP & NOT NULL & Report creation time & NOW() \\
    updated\_at & TIMESTAMP & NOT NULL & Last update timestamp & NOW() \\
    \bottomrule
  \end{tabularx}
\end{table}

\subsection{Activity Logs Table}

The \texttt{activity\_logs} table provides an audit trail of all significant user actions.

\begin{table}[H]
  \centering
  \caption{Activity Logs Table Schema}
  \label{tab:schema-activity}
  \begin{tabularx}{\textwidth}{@{}l l l X l@{}}
    \toprule
    \textbf{Column} & \textbf{Type} & \textbf{Constraints} & \textbf{Description} & \textbf{Default} \\
    \midrule
    id & INTEGER & PK, AUTO\_INCREMENT & Unique log entry ID & --- \\
    user\_id & INTEGER & FK $\to$ users.id, NOT NULL, INDEX & Action performer reference & --- \\
    action & VARCHAR(100) & NOT NULL & Action type identifier & --- \\
    details & TEXT & NULLABLE & Descriptive action details & NULL \\
    ip\_address & VARCHAR(45) & NULLABLE & Client IP address & NULL \\
    created\_at & TIMESTAMP & NOT NULL & Action timestamp & NOW() \\
    \bottomrule
  \end{tabularx}
\end{table}

\section{Relationships}

The SafeGuard database employs the following relationships:

\begin{table}[H]
  \centering
  \caption{Database Relationships}
  \label{tab:relationships}
  \begin{tabularx}{\textwidth}{@{}l l l X@{}}
    \toprule
    \textbf{Parent Table} & \textbf{Child Table} & \textbf{Type} & \textbf{Description} \\
    \midrule
    users & emergency\_contacts & One-to-Many & A user can have many emergency contacts; each contact belongs to one user. Cascade delete ensures contacts are removed when a user is deleted. \\
    users & sos\_alerts & One-to-Many & A user can trigger many SOS alerts; each alert is associated with one user. Cascade delete ensures alerts are removed with the user. \\
    users & incident\_reports & One-to-Many & A user can file many incident reports; each report belongs to one user. Cascade delete ensures reports are removed with the user. \\
    users & activity\_logs & One-to-Many & A user generates many activity log entries; each entry is associated with one user. Cascade delete ensures logs are removed with the user. \\
    \bottomrule
  \end{tabularx}
\end{table}

All foreign key constraints use \texttt{ON DELETE CASCADE}, ensuring referential integrity is maintained when parent records are removed. This design simplifies data management while ensuring no orphaned records exist in child tables.

\section{Primary and Foreign Keys}

\begin{itemize}
  \item \textbf{Primary Keys:} All tables use auto-incrementing integer primary keys (\texttt{id}), except the \texttt{incident\_id} field in \texttt{incident\_reports} which serves as a secondary unique key with the format \texttt{INC-XXXXXXXX}.

  \item \textbf{Foreign Keys:} The \texttt{user\_id} column in \texttt{emergency\_contacts}, \texttt{sos\_alerts}, \texttt{incident\_reports}, and \texttt{activity\_logs} tables references the \texttt{id} column in the \texttt{users} table, establishing the owner relationship.

  \item \textbf{Indexes:} Indexes are created on all foreign key columns for query optimization, as well as on frequently queried columns including \texttt{email} (users), \texttt{incident\_id} (incident\_reports), and \texttt{status} (sos\_alerts).
\end{itemize}

\section{Data Dictionary}

\begin{table}[H]
  \centering
  \caption{Data Dictionary: ENUM Types}
  \label{tab:data-dict}
  \begin{tabularx}{\textwidth}{@{}l l X@{}}
    \toprule
    \textbf{ENUM Field} & \textbf{Values} & \textbf{Description} \\
    \midrule
    user.role & user, admin & System role: standard user or administrator \\
    emergency\_contacts.priority & low, medium, high, critical & Contact priority for emergency notification \\
    sos\_alerts.alert\_type & manual, automatic, panic, medical, fire, security & Classification of the emergency \\
    sos\_alerts.status & active, acknowledged, resolved, cancelled & Alert lifecycle state \\
    incident\_reports.incident\_type & theft, assault, accident, fire, harassment, vandalism, medical, other & Nature of the reported incident \\
    incident\_reports.severity & low, medium, high, critical & Impact level of the incident \\
    incident\_reports.status & pending, investigating, resolved, closed & Investigation progress state \\
    \bottomrule
  \end{tabularx}
\end{table}

\begin{table}[H]
  \centering
  \caption{Data Dictionary: Data Types}
  \label{tab:data-types}
  \begin{tabularx}{\textwidth}{@{}l l X@{}}
    \toprule
    \textbf{Data Type} & \textbf{SQL Equivalent} & \textbf{Usage} \\
    \midrule
    VARCHAR(n) & VARCHAR(n) & Variable-length string with specified maximum length \\
    TEXT & TEXT & Variable-length string for long content (descriptions, addresses) \\
    INTEGER & INT & Whole numbers (IDs, counts) \\
    BOOLEAN & TINYINT(1) / BIT & True/false values \\
    NUMERIC(10,8) & DECIMAL(10,8) & GPS coordinates with 8 decimal places of precision \\
    TIMESTAMP & DATETIME & Date and time values \\
    ENUM & ENUM & Predefined set of allowed values \\
    \bottomrule
  \end{tabularx}
\end{table}

\section{Sample Data}

This section presents sample data stored in each database table to illustrate the actual data patterns and relationships in the SafeGuard system.

\subsection{Users Table Data}

Figure~\ref{fig:data-users} shows sample records from the \texttt{users} table, demonstrating user account information including credentials, profile data, and role assignments.

\begin{figure}[H]
  \centering
  \includegraphics[width=0.95\textwidth]{db_users.png}
  \caption{Sample Data: Users Table}
  \label{fig:data-users}
\end{figure}

\subsection{Emergency Contacts Table Data}

Figure~\ref{fig:data-contacts} shows sample emergency contact records with priority classifications.

\begin{figure}[H]
  \centering
  \includegraphics[width=0.95\textwidth]{db_emergency_contacts.png}
  \caption{Sample Data: Emergency Contacts Table}
  \label{fig:data-contacts}
\end{figure}

\subsection{SOS Alerts Table Data}

Figure~\ref{fig:data-sos} shows sample SOS alert records with geolocation and status tracking.

\begin{figure}[H]
  \centering
  \includegraphics[width=0.95\textwidth]{db_sos_alerts.png}
  \caption{Sample Data: SOS Alerts Table}
  \label{fig:data-sos}
\end{figure}

\subsection{Incident Reports Table Data}

Figure~\ref{fig:data-incidents} shows sample incident reports with type, severity, and status tracking.

\begin{figure}[H]
  \centering
  \includegraphics[width=0.95\textwidth]{db_incident_reports.png}
  \caption{Sample Data: Incident Reports Table}
  \label{fig:data-incidents}
\end{figure}

\subsection{Activity Logs Table Data}

Figure~\ref{fig:data-activity} shows sample activity log entries tracking user actions.

\begin{figure}[H]
  \centering
  \includegraphics[width=0.95\textwidth]{db_activity_logs.png}
  \caption{Sample Data: Activity Logs Table}
  \label{fig:data-activity}
\end{figure}
